\chapter{Introducción}

\section{Contexto clínico y motivación}

Las masas tumorales intracraneales constituyen un grupo de lesiones de elevada
relevancia clínica: su localización dentro del cráneo hace que incluso lesiones
histológicamente poco agresivas puedan producir síntomas por efecto masa, edema o
compresión de estructuras críticas, y su hallazgo desencadena con frecuencia una
cadena de decisiones diagnósticas y terapéuticas de gran impacto. La resonancia
magnética (RM) cerebral es la técnica de referencia para su detección y
caracterización, por lo que la rapidez con la que un estudio con sospecha de masa
llega a ser revisado puede condicionar el itinerario del paciente.

En paralelo, el volumen de exploraciones de imagen ha crecido de forma sostenida,
y con él la carga de trabajo de los servicios de radiología. En contextos de alta
demanda ---guardias, turnos nocturnos o acumulación de listas de lectura--- el
orden en que se revisan los estudios deja de ser indiferente. De ahí el interés
por los sistemas de priorización automática, que reordenan la lista de trabajo
para que los estudios potencialmente urgentes sean atendidos antes, sin sustituir
el juicio del especialista~\cite{titano2018cranial}. La priorización de estudios
de RM cerebral con sospecha de masa tumoral es, en este sentido, una tarea con
interés asistencial directo y, a la vez, metodológicamente acotada: no exige
emitir un diagnóstico completo, sino estimar si un estudio merece revisión
prioritaria.

\section{Planteamiento del problema}

Formular esa priorización como una tarea de clasificación binaria
tumor/no-tumor a nivel de estudio, y abordarla con aprendizaje profundo a partir
de datos públicos, es un enfoque atractivo y habitual. Sin embargo, los datasets
públicos de RM cerebral suelen ser \emph{monoclase}: unos aportan únicamente
casos con tumor y otros únicamente sujetos sin lesión. Al combinarlos para
construir el conjunto de entrenamiento, se corre el riesgo de que la clase quede
correlacionada con la procedencia del estudio, es decir, con el centro, el equipo
y el protocolo de adquisición.

Cuando esa correlación es alta, una métrica de validación interna elevada deja de
ser informativa: un modelo puede alcanzarla reconociendo la firma de adquisición
de cada cohorte en lugar de la señal patológica. Este fenómeno ---el aprendizaje
de atajos espurios en lugar de la característica de interés--- está bien
documentado en aprendizaje profundo~\cite{geirhos2020shortcut} y, de forma
específica, en imagen médica, donde modelos aparentemente exitosos resultaron
apoyarse en marcadores de procedencia antes que en signos de
enfermedad~\cite{zech2018pneumonia}. En consecuencia, un buen resultado interno es
una condición necesaria pero no suficiente para afirmar que un sistema detecta
lesión, y exige una verificación específica antes de ser interpretado como tal.

\section{De la hipótesis clínica inicial a la auditoría metodológica del modelo}

Este trabajo parte de una pregunta clínica legítima ---si una red convolucional
3D puede priorizar estudios de RM con sospecha de masa tumoral--- y la aborda de
forma deliberadamente crítica. En lugar de dar por buena la métrica que arroja un
clasificador entrenado sobre el conjunto multi-fuente, adopta como objeto de
estudio la propia \emph{validez} de ese rendimiento: la pregunta operativa se
desplaza de «¿qué rendimiento se obtiene?» a «¿ese rendimiento refleja detección
de lesión o reconocimiento de la procedencia del estudio?».

Este desplazamiento no es una renuncia a la motivación clínica, sino la postura
analítica que el problema exige cuando clase y dominio pueden confundirse. El
título de la memoria lo encapsula: el núcleo no es entregar un clasificador, sino
realizar un \emph{análisis crítico del sesgo de dominio} en este tipo de sistemas.
Bajo esta perspectiva, el clasificador deja de ser el producto y pasa a ser el
sujeto de una auditoría: un caso de estudio sobre el que medir, de forma
reproducible, si el aprendizaje es genuino o espurio.

\section{Objetivos y contribuciones}

El objetivo de este trabajo es auditar, sobre un clasificador tumor/no-tumor
basado en una CNN 3D y entrenado con RM cerebral pública multi-fuente, si su
rendimiento aparente refleja la detección de lesión o la explotación de un
\emph{confound} dominio$\leftrightarrow$clase, y proponer un protocolo
reproducible para diagnosticar este modo de fallo. La hipótesis de trabajo es que,
cuando la etiqueta coincide con el dataset de origen, el rendimiento interno es
atribuible al confound y no constituye evidencia de capacidad de detección; el
trabajo se diseña para contrastar esa hipótesis de forma sistemática.

De este planteamiento se derivan tres contribuciones. La primera es un
\textbf{pipeline 3D reproducible} para clasificación de RM cerebral multi-fuente,
con preprocesado homogéneo, partición a nivel de sujeto y semillas fijas. La
segunda, y principal, es un \textbf{protocolo de auditoría anti-confounding} que
encadena un baseline trivial de intensidad, validación cruzada por dominio
(\emph{leave-one-dataset-out}), validación intra-dominio y análisis del espacio
latente, junto con el descarte explícito de fugas por partición y duplicados. La
tercera es la \textbf{caracterización del modo de fallo}: la identificación del
confound dominio$\leftrightarrow$clase como causa del rendimiento aparente y de
las recomendaciones metodológicas que de ello se siguen para el trabajo con datos
públicos heterogéneos. La formulación detallada de estos objetivos se presenta en
el capítulo siguiente.

\section{Estructura de la memoria}

La memoria se organiza como sigue. El presente capítulo (1) introduce el problema
y los objetivos, que se enuncian de forma detallada en el capítulo~2. El
capítulo~3 revisa el estado del arte ---contexto clínico, RM cerebral, aprendizaje
profundo en imagen médica, sesgo de dominio y triaje--- y sitúa el hueco que
aborda el trabajo. El capítulo~4 describe los materiales y datos, y el capítulo~5
la metodología (modelo, entrenamiento y partición). El capítulo~6 detalla el
diseño experimental y el protocolo de validación, sin cifras, y el capítulo~7
reporta los resultados de cada experimento. El capítulo~8 los discute e
interpreta a la luz de la literatura. Finalmente, el capítulo~9 recoge las
conclusiones, el capítulo~10 las limitaciones y el trabajo futuro, y los anexos
reúnen el material de apoyo.
